\chapter{التعرّف على الصور والفيديو}
\chaptermark{التعرّف}
\label{sec:recognition}
\begin{chaptergoals}
\item لماذا يجب أن يتجاهل التعرّف الإزاحة والمقياس والإضاءة، ولماذا تفشل المقارنة بكسلًا ببكسل؛
\item كيف يتعرّف خط أنابيب يناوب بين ”و“ على الأجزاء و”أو“ على المواضع في $150\ms$؛
\item كيف تتعلم قاعدة هيب مع أثر اللاتغيّر من فيديو بلا أسماء؛
\item بمَ يختلف الدماغ عن الشبكة الالتفافية، وماذا تكشف الخدع البصرية عنه.
\end{chaptergoals}

\section{المسألة: الشيء نفسه ببكسلات مختلفة}

أوصل الفصل~\ref{sec:sensors} الصورة إلى مدخل الآلة: نحو مليون عصبون مخرج في العين، وتدفق مضغوط لا يُنقل فيه في الغالب إلا الحواف والتغيرات. لكن بين المستشعر والفعل خطوة سكتنا عنها حتى الآن. القطة في الصورة ليست مجموعة من البكسلات. أزِحها نصف إطار، أو أدِرها، أو أبعِدها، أو أطفئ نصف الضوء، فتتغير كل البكسلات تقريبًا، ومع ذلك يجب أن يبقى الجواب ”قطة“ كما هو. يسمي المهندس هذا \emph{اللاتغيّر}: يجب ألا يتغير جواب المصنِّف مع الإزاحة والمقياس والدوران والإضاءة.

تظهر الصعوبة جليًا في تجربة بسيطة. لنأخذ ”شبكية“ أحادية البعد من $24$ نقطة وشكلين من خمس نقاط يمكن أن يقعا في أي مكان. المصنِّف الذي يقارن الدخل بنموذج محفوظ في موضع واحد يصيب في $49\text{--}52\%$ من الحالات، أي بقدر رمي عملة. فالإزاحة تحطم المقارنة بكسلًا ببكسل تمامًا. لا بد من مخطط آخر.

\section{خط أنابيب من مراحل}

نعرف من الفصل~\ref{sec:code} مدى سرعة الدماغ في ذلك: يُتعرَّف على حيوان في صورة معروضة لحظيًا في نحو $150\ms$، وتمر الإشارة بنحو عشر مراحل، من $10$ إلى $20\ms$ لكل منها~\cite{Thorpe1996}. وفي كل مرحلة لا يلحق العصبون أن يُصدر إلا صفرًا من النبضات أو نبضة واحدة. إذن التعرّف السريع مرور واحد عبر خط الأنابيب، بلا تفكير طويل ولا تكرار. والسؤال: ماذا تفعل كل مرحلة؟

الجواب الذي أوحت به القياسات في المراحل الأولى للبصر~\cite{HubelWiesel1962} يتألف من عمليتين متناوبتين (الشكل~\ref{fig:conveyor}).

\emph{الانتقائية.} ينظر عصبون المرحلة $S$ إلى جزء صغير من الدخل، ولا يطلق إلا إذا وُجد فيه تركيب معيّن من الأجزاء: هذه الحافة، وبجوارها تلك. هذا ”و“ على الأجزاء، وهو العصبون العتبي من الفصل~\ref{sec:glutcomputer} الذي عتبته أعلى مما يعطيه أي جزء وحده.

\emph{اللاتغيّر.} يجمع عصبون المرحلة $C$ مخارج عصبونات $S$ كثيرة تبحث عن التركيب نفسه في مواضع مختلفة، ويطلق إذا وُجد التركيب في أي مكان. هذا ”أو“ على المواضع، أو بتعبير أدق: اختيار الأكبر.

وفي النموذج الأحادي البعد تُكتب العمليتان هكذا:
\begin{empheq}[box=\keybox]{equation}
s_{f,p} \;=\; \Bigl[\sum_i t_{f,i}\,x_{p+i} - \theta\Bigr]_+ ,
\qquad
c_f \;=\; \max_p\, s_{f,p} ,
\label{eq:scstage}
\end{empheq}
حيث $x$ الدخل، و$t_{f,i}$ نموذج السمة $f$، و$p$ الموضع، و$\theta$ العتبة. التعبير الأول يجعل الجواب انتقائيًا، والثاني يجعله مستقلًا عن الموضع. وتحصل الشبكة على أكبر المداخل بأدوات الفصل~\ref{sec:gaba}: التثبيط المتبادل يُبقي أقوى الدخل (القضية~\ref{prop:wta})، والقسمة على النشاط الكلي تسوّي الاستجابات عند سطوع مختلف~\cite{Carandini2012}.

\begin{figure}[t]
\centering
\begin{tikzpicture}[>=Stealth,
  px/.style={draw,minimum size=0.36cm,inner sep=0pt},
  on/.style={px,fill=blue!55!black},
  snode/.style={circle,draw,very thick,minimum size=0.62cm,inner sep=0pt,font=\scriptsize,fill=blue!6},
  cnode/.style={circle,draw,very thick,minimum size=0.8cm,inner sep=0pt,font=\scriptsize,fill=red!10}]
% two inputs: the same shape at two positions
\foreach \row/\sh/\lab in {0/1/{الإطار 1},1/7/{الإطار 2}}{
  \begin{scope}[yshift=-\row*1.0cm]
  \node[left,font=\scriptsize] at (-0.25,0) {\lab};
  \foreach \i in {0,...,13}{\node[px] at (\i*0.4,0) {};}
  \foreach \d in {0,1,3,4}{\node[on] at ({(\sh+\d)*0.4},0) {};}
  \end{scope}}
% S stage: detectors of the shape at several positions
\foreach \k in {0,...,5}{
  \node[snode] (S\k) at ({0.8+0.8*\k},1.7) {$S$};}
\node[font=\scriptsize,left] at (-0.25,1.7) {المرحلة $S$: ”و“ على الأجزاء};
\draw[->,thick,blue!60!black] (1.2,0.25) -- (S1.south);
\draw[->,thick,orange!85!black] (3.6,0.25) -- (S4.south);
\node[font=\scriptsize,blue!60!black,anchor=east] at (1.25,0.85) {الإطار 1};
\node[font=\scriptsize,orange!85!black,anchor=west] at (3.9,0.85) {الإطار 2};
% C stage
\node[cnode] (C) at (2.8,3.25) {$C$};
\node[font=\scriptsize,left] at (-0.25,3.25) {المرحلة $C$: ”أو“ على المواضع};
\foreach \k in {0,...,5}{\draw[->,gray!60!black] (S\k.north) -- (C);}
\draw[->,very thick,red!70!black] (C.north) -- ++(0,0.6) node[above,font=\scriptsize,black] {”الشكل موجود“ في الإطارين كليهما};
% next stages
\draw[decorate,decoration={brace,amplitude=5pt},gray!60!black] (6.3,1.4) -- (6.3,-1.3);
\node[right,font=\scriptsize,align=left] at (6.5,0.05) {الشيء نفسه،\\وبكسلات\\أخرى};
\draw[->,thick,densely dashed,gray!60!black] (3.6,3.25) -- (5.9,3.25) node[right,font=\scriptsize,black,align=left] {الأزواج التالية\\$S$، $C$: أكبر،\\وأعقد، وأكثر لاتغيّرًا};
\end{tikzpicture}
\caption{زوج واحد من مراحل خط الأنابيب. تبحث عصبونات $S$ عن تركيب الأجزاء نفسه في مواضع مختلفة؛ ويستجيب العصبون $C$ إذا وُجد التركيب في أي مكان~\eqref{eq:scstage}. الشكل في الإطار 1 وفي الإطار 2 يثير عصبونات $S$ مختلفة، لكنه يثير العصبون $C$ نفسه. والأزواج التالية من المراحل تكرر الشيء نفسه على تراكيب أكبر فأكبر وأعقد فأعقد.}
\label{fig:conveyor}
\end{figure}

في النموذج الأحادي البعد نفسه يتعرف زوج المراحل~\eqref{eq:scstage} على الشكل في أي مكان بلا أخطاء، وإذا قُلبت كل نقطة من الدخل عشوائيًا باحتمال $0.1$ فإنه يصيب في $86\%$ من الحالات، وباحتمال $0.2$ في $70\%$. أما رمي العملة فيعطي النصف كما من قبل. ضع عدة أزواج كهذه بعضها فوق بعض: الأول يبحث عن الحواف، والتالي عن تراكيب الحواف، وأعلى منه عن أجزاء الأشياء، وفي القمة عن الأشياء كاملة. ومع كل زوج تكبر التراكيب، ويقل اعتماد الجواب على مكان الشيء وهيئته~\cite{Riesenhuber1999}.

\begin{keyblock}
\begin{proposition}[خط أنابيب الانتقائية واللاتغيّر]
\label{prop:conveyor}
التعرّف السريع مرور واحد عبر نحو عشر مراحل. في كل زوج من المراحل يجعل ”و“ على الأجزاء~\eqref{eq:scstage} الجواب انتقائيًا، ويجعله ”أو“ على المواضع مستقلًا عن الإزاحة. وتناوب هاتين العمليتين يعطي جوابًا يميّز الأشياء ولا يكاد يتوقف على مكانها وهيئة رؤيتها. بنية المراحل الأولى وسرعة المرور مقيستان؛ أما أن السلسلة كلها تُختزل إلى هاتين العمليتين فتبسيط.
\end{proposition}
\end{keyblock}

إذا بدت هذه البنية مألوفة، فهذا طبيعي. فتناوبها هذا بالذات، أي البحث عن السمة في كل المواضع ثم الدمج على المواضع المتجاورة، هو ما نقله المهندسون إلى الشبكات الاصطناعية \emph{الالتفافية}~\cite{Fukushima1980}. ثم ساروا مؤخرًا في الاتجاه المعاكس: تبيّن أن الشبكات الالتفافية المدرَّبة على التعرف على الأشياء أفضل نموذج معروف لاستجابات العصبونات في المراحل العليا للبصر~\cite{Yamins2014}. والنتيجة في القمة تُقرأ ببساطة: من الترددات الكلية لنحو مئة عصبون في المرحلة الأخيرة تتعرف وحدة قرار خطية على الشيء بعد عُشر ثانية من عرضه~\cite{Hung2005}. هذه شفرة التردد الجماعية من الفصل~\ref{sec:code}.

\section{المعلّم الغائب: التعلم من الفيديو}

تُدرَّب الشبكة الالتفافية على ملايين الصور المعنونة بأسمائها. أما الدماغ فلا يكاد أحد يعطيه أسماء: يرى الطفل الأشياء ساعات، ولا يسمع كلمة ”فنجان“ إلا نادرًا. فمن أين تعرف مراحل $C$ أيّ عصبونات $S$ تدمج؟ فلكي ندمج ”هذا الشكل هنا“ مع ”هذا الشكل هناك“ يجب أن نعرف مسبقًا أنه الشكل نفسه.

والتلميح يأتي من الفيديو نفسه. العالم يتغير باستمرار: الشيء في مجال الرؤية يبقى الشيء نفسه وهو يتزحزح ويدور وتتغير إضاءته، ولا ينتقل البصر إلى غيره إلا أحيانًا. فكل ما يتغير \emph{ببطء} يرجح أنه يخص الشيء، وكل ما يتغير بسرعة يخص موضعه وهيئته~\cite{WiskottSejnowski2002}. إذن يكفي العصبون $C$ قاعدة هيب من الفصل~\ref{sec:dofa} مع أثر: أن يربط ما جاء متتابعًا، لا ما جاء متزامنًا فقط~\cite{Foldiak1991}:
\begin{empheq}[box=\keybox]{equation}
\tau_{\text{tr}}\,\frac{\dd\bar y_k}{\dd t} \;=\; -\bar y_k + y_k ,
\qquad
\frac{\dd\mathbf w_k}{\dd t} \;=\; \eta\,\bar y_k\,\bigl(\mathbf x - \mathbf w_k\bigr) .
\label{eq:traceinvariance}
\end{empheq}
هنا $y_k$ نشاط العصبون $C$ رقم $k$ الذي فاز في التنافس عبر التثبيط، و$\bar y_k$ أثره، وهو دارة \foreignlanguage{english}{RC} نفسها التي في القاعدة~\eqref{eq:threefactor}، و$\mathbf x$ مخارج عصبونات $S$. العصبون الذي فاز على الهيئة الأولى للشيء لا يزال يتذكر ذلك حين تأتي الهيئة الثانية، فيجذبها إليه هي أيضًا.

\begin{figure}[t]
\centering
\begin{tikzpicture}[>=Stealth,x=1.0cm,y=1.0cm]
\foreach \i/\lab in {0/1,1/2,2/3,3/4}{
  \node[draw,thick,fill=blue!8,minimum width=0.9cm,minimum height=0.55cm,font=\scriptsize] at (\i+0.5,2.2) {$A$ في \lab};}
\foreach \i/\lab in {0/4,1/3,2/2,3/1}{
  \node[draw,thick,fill=orange!12,minimum width=0.9cm,minimum height=0.55cm,font=\scriptsize] at (\i+6.5,2.2) {$B$ في \lab};}
\node[font=\scriptsize,gray!40!black] at (5.0,2.2) {توقف};
\draw[->,gray!70!black] (0,0) -- (10.4,0) node[below left,font=\scriptsize,black] {الزمن};
% traces
\draw[thick,blue!60!black] (0,0.05) .. controls (0.4,1.2) and (0.8,1.3) .. (1.2,1.35) -- (4,1.4) .. controls (4.6,0.5) and (5.2,0.15) .. (6,0.08) -- (10,0.05);
\draw[thick,orange!85!black] (0,0.05) -- (6,0.05) .. controls (6.4,1.2) and (6.8,1.3) .. (7.2,1.35) -- (10,1.4);
\node[font=\scriptsize,blue!60!black,anchor=west] at (1.4,1.65) {الأثر $\bar y_1$ يدوم طوال مرور $A$};
\node[font=\scriptsize,orange!85!black,anchor=west] at (7.2,1.65) {الأثر $\bar y_2$};
\draw[decorate,decoration={brace,amplitude=4pt},gray!60!black] (4.0,2.6) -- (0,2.6);
\draw[decorate,decoration={brace,amplitude=4pt,mirror},gray!60!black] (0,-0.35) -- (4,-0.35) node[midway,below=4pt,font=\scriptsize,black] {كل هيئات $A$ ترتبط بالعصبون 1};
\end{tikzpicture}
\caption{كيف يعلّم الفيديو اللاتغيّر، مخطط. يمر الشيء $A$ بأربعة مواضع، ثم توقف، ثم الشيء $B$. أثر~\eqref{eq:traceinvariance} العصبون الذي فاز على الهيئة الأولى لـ$A$ يدوم طوال المرور، فترتبط كل هيئات $A$ بعصبون واحد. والتوقف يمحو الأثر، فيذهب $B$ إلى عصبون آخر.}
\label{fig:tracevideo}
\end{figure}

اختبرنا ذلك على نموذج (الشكل~\ref{fig:tracevideo}). عصبونا $C$ يتنافسان على مداخل من $16$ عصبون $S$: شكلان في ثمانية مواضع. يمر الشكل عبر كل المواضع، $50\ms$ لكل موضع، ثم توقف $0.3\,\text{s}$، ثم الشكل العشوائي التالي. بلا أي أسماء. إذا كان الأثر قصيرًا، $\tau_{\text{tr}}=5\ms$، اقتسم العصبونان الدخل بحسب \emph{الموضع}، ولم يطابق الجوابُ الشكلَ بأفضل من رمي عملة: $52\%$ في المتوسط على عشرة تشغيلات. وإذا كان الأثر مماثلًا لزمن عرض موضع واحد، ابتداءً من $\tau_{\text{tr}}\approx20\ms$، جمع كل عصبون \emph{كل} مواضع شكله (عند $20$ و$50$ و$200\ms$ في التشغيلات العشرة كلها): تعلّم اللاتغيّر من الفيديو وحده. والتوقف بين الأشياء مهم: بدونه ينساب الأثر إلى الشيء التالي فيخلط بينهما.

والشيء نفسه يُرى في التجربة. إذا استُبدل الشيء خفيةً في التجربة بينما تقفز العين من مكان إلى آخر، فإن عصبونات المراحل العليا تبدأ في غضون ساعة أو ساعتين بالاستجابة للشيء البديل كأنه الشيء نفسه: إنها تربط ما جاء متتابعًا~\cite{LiDiCarlo2008}. فاللاتغيّر في الدماغ يُتعلَّم فعلًا من استمرارية الزمن. أما إلى أي حد تفسر هذه القاعدة كل تعلم البصر، ففرضية.

\section{الفيديو حركة}

ليس الفيديو مجرد تدفق من الإطارات للتعلم، بل هو المسألة نفسها أيضًا: ما الذي يتحرك، وإلى أين، وبأي سرعة. الخطوة الأولى معروفة لنا: كاشف الاتجاه من الفصل~\ref{sec:gaba} (الشكل~\ref{fig:direction})، الذي يُطفئ فيه التثبيطُ المتأخر الحركةَ في اتجاه واحد. تنتشر هذه الكواشف في كل مجال الرؤية، ثم يُعامَل معها كما يُعامَل مع سمات الشكل: المراحل التي تجمع كواشف كثيرة للاتجاه نفسه في مواضع مختلفة تستجيب لحركة شيء كبير أينما كان. هذا زوج ”و“ و”أو“ نفسه~\eqref{eq:scstage}، إلا أن السمة ليست ”حافة“، بل ”حافة انزاحت خلال زمن $d$“.

والفيديو فوق ذلك قابل للتنبؤ: الإطار التالي يكاد يطابق السابق. بحسب فرضية \emph{الترميز التنبؤي} ترسل المراحل العليا إلى الدنيا تنبؤًا، ولا يصعد إلى الأعلى في الغالب إلا الخطأ، أي ما لم يحسب التنبؤ حسابه~\cite{RaoBallard1999}. هذا هو الاقتصاد نفسه في النبضات الذي في القضية~\ref{prop:economy}: ادفع ثمن الجديد، لا ثمن المعروف أصلًا. بعض المشاهدات تتفق مع هذه الفرضية، لكنها لم تُثبت بوصفها البنية العامة لخط الأنابيب.

\section{الدماغ والشبكة الالتفافية}

إلى أي مدى يذهب التشابه؟ في التعرّف السريع على شيء واحد هو عميق حقًا~\cite{Yamins2014}. لكن في الدماغ ما ليس في الشبكة الالتفافية البسيطة.

\emph{الوصلات الراجعة.} خط أنابيب الدماغ ليس أماميًا فقط: لكل مرحلة وصلات إلى الخلف وإلى الجانب. في الصور الصعبة، كالمتداخلة أو ضعيفة الإضاءة، تتكوّن استجابة المراحل العليا متأخرة، وهذه الاستجابات المتأخرة تصفها الشبكات ذات الوصلات الراجعة أفضل من الشبكات الأمامية~\cite{Kar2019}. المرور الواحد يحل الحالات السهلة، والصعبة تتطلب عدة دورات.

\emph{المتانة.} يمكن خداع الشبكة الاصطناعية بتزييف في البكسلات لا تلحظه العين~\cite{Szegedy2014}: تغيير لا يراه الإنسان يحوّل ”الباندا“ إلى ”جِبّون“~\cite{Goodfellow2015}. والبصر البشري أمتن بكثير أمام هذا التزييف، لكنه ليس محصنًا: الصور التي تخدع شبكات كثيرة معًا تزيح اختيار الإنسان أيضًا عند العرض القصير~\cite{Elsayed2018}. لماذا تختلف المتانة إلى هذا الحد؟ سؤال مفتوح؛ والمرشحون: الضجيج، والقسمة على النشاط الكلي، والوصلات الراجعة.

\emph{المعلّم.} تحتاج الشبكة إلى ملايين الأمثلة المعنونة، أما الدماغ فيحتاج في الغالب إلى فيديو بلا أسماء~\eqref{eq:traceinvariance} وقليل من التلميحات من \nt{DOPA} عما هو مهم.

\emph{الطاقة.} الدماغ كله نحو $20$ واطًا (القضية~\ref{prop:economy})، والتعرّف لا يشغل إلا جزءًا من ذلك؛ أما الشبكة الالتفافية المدرَّبة على مسرّع رسوميات فتستهلك مئات الواطات.

\section{الخداع البصري: أين تخطئ الآلة ولماذا}
\label{sec:illusions}

المهندس الذي يريد أن يفهم دارة غيره يغذيها بإشارات غير مألوفة، وينظر أين تخطئ. وللبصر إشارات كهذه ابتُكرت منذ زمن بعيد، هي الخدع البصرية. الخدعة ليست عطلًا. كل خدعة تشير إلى آلية بعينها في الآلة، تعمل صحيحًا في الظروف العادية، وتكشف نفسها في صورة مختارة خصيصًا.

\emph{توقعات معقولة.} الدخل مشوَّش بالضجيج، والآلة تكمله بما يحدث عادةً في العالم. الحركات البطيئة أكثر من السريعة؛ وحين يكون الدخل غير موثوق، كأن يكون تباين الصورة ضعيفًا، ينزاح تقدير السرعة نحو البطء، فيبدو الشيء الباهت أبطأ من الساطع~\cite{Weiss2002}. وبالطريقة نفسها تُفسَّر خدع كثيرة في السطوع: نحن لا نرى سطوع البكسل، بل السطح الذي طابقه هذا السطوع في أغلب الأحيان في الماضي~\cite{Purves2011}. وصورة الفستان الذي يراه بعضهم أبيض وذهبيًا ويراه آخرون أزرق وأسود هي الآلية نفسها: الصورة ملتبسة، ويختلف الناس في الإضاءة التي يفترضونها دون وعي~\cite{LaferSousa2015,Wallisch2017}. الفروق بين الناس مقيسة؛ أما أن الدماغ يجمع هنا الدخل والتوقع وفق قواعد نظرية الاحتمالات ففرضية راسخة الأسس.

\emph{ضبط الكسب.} انظر دقيقة إلى شلال، ثم إلى صخور ساكنة: ستبدو الصخور كأنها تطفو إلى أعلى. قناتا ”إلى أسفل“ و”إلى أعلى“ زوج أسلاك كما في~\eqref{eq:dualrail}؛ وقد خفّض ضبط الكسب~\eqref{eq:nakarushton} القناة المتعبة، فاختل توازن الزوج. وخدع الميل والتباين التي يغيّر فيها المحيطُ مظهرَ المركز تُفسَّر بالقسمة على نشاط الجيران~\cite{Schwartz2009}، كما في الفصل~\ref{sec:gaba}. وهناك مثال مفيد على الحذر: البقع الداكنة عند تقاطعات الخطوط البيضاء في الشبكة المربعة فُسرت طويلًا بتثبيط الجيران البسيط، لكن التجارب على شبكة مربعة معدّلة أظهرت أن هذا التفسير لا يصح~\cite{SchillerCarvey2005}.

\emph{تعويض التأخيرات.} يستغرق المسار من العين إلى المراحل العليا عشرات الملّي ثانية، والشيء المتحرك يلحق أن ينزاح خلالها. إذا ومضت نقطة ساكنة بجانبه، بدا الشيء متقدمًا على الومضة مع أنهما كانا متطابقين~\cite{Nijhawan1994}. وبحسب أحد التفسيرات تمدّ الآلة الحركة إلى الأمام لتعوّض التأخير، كما في الفصل~\ref{sec:muscles}: مع تغذية راجعة متأخرة لا بد من التنبؤ. ولهذه الخدعة تفسيرات عدة، والخلاف لم يُحسم.

\emph{خطأ في شفرة التوقيت.} صورة ساكنة من حلقات ذات انتقالات لونية حادة تبدو كأنها تدور. فالمناطق عالية التباين تُعالَج أسرع من الباهتة، وتقرأ كواشفُ الاتجاه فرقَ التأخيرات~\eqref{eq:latency} على أنه حركة~\cite{Conway2005}. ويطلق الخدعة قفزات العين الصغيرة اللاإرادية والرمش: كل قفزة تجدد الدخل، فترى الكواشف ”حركة“ من جديد~\cite{OteroMillan2012}. هذه شفرة التوقيت من الفصل~\ref{sec:code} وقد قُرئت خطأً.

\emph{صورتان في صورة.} هيكل مكعب يواجهنا تارة بوجه وتارة بآخر، أو صورتان مختلفتان في العينين: يقفز الإدراك كل بضع ثوانٍ. أفضل النماذج تثبيط متبادل بين التأويلين، وتعب بطيء، وضجيج~\cite{MorenoBote2007}، أي المخطط نفسه~\eqref{eq:halfcentre} الذي ولّد في الفصل~\ref{sec:muscles} إيقاع الخطو. يحدد زمنَ التبدل الضجيجُ والتعب؛ وبحسب النموذج~\cite{MorenoBote2007} الدور الأكبر للضجيج، والتعب يساعد فقط.

وماذا عن الشبكات الاصطناعية؟ الشبكة المدرَّبة على التنبؤ بالإطار التالي من الفيديو ترى الدوران في الحلقات الساكنة نفسها، ولا تراه في الصور الضابطة~\cite{Watanabe2018}. والشبكات المدرَّبة على تنقية الصور من الضجيج وزيادة حدتها تنخدع بخدع اللون والسطوع نفسها التي ينخدع بها الناس~\cite{GomezVilla2020}. إذن بعض الخدع نتيجة للمسألة نفسها التي تتعلمها الآلة، لا لخصائص النسيج العصبي. وأي الخدع يختص بها الدماغ وحده؟ هذا اختبار جيد لأي نموذج للبصر.

\begin{keyblock}
\begin{proposition}[الخدع بصمات للآليات]
\label{prop:illusions}
تنشأ الخدعة حين تتلقى آلية نافعة في العالم العادي دخلًا غير مألوف: التوقع يغلب الدخل غير الموثوق، وضبط الكسب يزيح زوج القنوات، وتعويض التأخير يقدّم الشيء إلى الأمام، وفرق التأخيرات يُقرأ حركةً، والتثبيط المتبادل مع الضجيج والتعب يقفز. كل خدعة إشارة اختبار تشير إلى آليتها. الخدع نفسها مقيسة؛ وتفسيراتها تتراوح بين الراسخ الأسس والمختلف فيه.
\end{proposition}
\end{keyblock}

\section{الخلاصة}

\begin{table}[t]
\centering
\footnotesize
\setlength{\tabcolsep}{4pt}
\renewcommand{\arraystretch}{1.15}
\begin{tabular}{>{\raggedright}p{3.0cm}>{\raggedright}p{4.6cm}>{\raggedright\arraybackslash}p{5.2cm}}
\toprule
ماذا تفعل الآلة & كيف & ما هو معروف \\
\midrule
تتعرف في $150\ms$ & مرور واحد عبر نحو عشر مراحل & السرعة مقيسة \\
تجعل الجواب انتقائيًا & ”و“ على الأجزاء، المرحلة $S$~\eqref{eq:scstage} & مقيس في المراحل الأولى \\
تجعل الجواب لا متغيرًا & ”أو“ على المواضع، المرحلة $C$؛ التثبيط والقسمة & مقيس في المراحل الأولى؛ وفي العليا تبسيط \\
تُصدر الجواب & ترددات نحو مئة عصبون في المرحلة الأخيرة، قراءة خطية & مقيس \\
تتعلم بلا أسماء & قاعدة هيب مع أثر~\eqref{eq:traceinvariance} على فيديو متصل & استبدال الشيء يغيّر الاستجابات، وهذا مقيس؛ وكونها القاعدة الرئيسية فرضية \\
ترى الحركة & كواشف الاتجاه، ثم زوج ”و“/”أو“ نفسه & مقيس \\
تقتصد في المتوقَّع & يصعد خطأ التنبؤ إلى الأعلى & فرضية \\
تحل الحالات الصعبة & وصلات راجعة، عدة دورات & الاستجابات المتأخرة ودور الوصلات الراجعة مقيسة \\
تخطئ على نحو متوقَّع & الخدع: التوقعات، ضبط الكسب، التأخيرات، التثبيط المتبادل مع الضجيج والتعب & الخدع مقيسة؛ والتفسيرات بين الراسخ والمختلف فيه \\
\bottomrule
\end{tabular}
\caption{كيف تتعرف الآلة على الصور والفيديو.}
\label{tab:recognition}
\end{table}

التعرّف مجمَّع من قطع كانت لدينا من قبل (الجدول~\ref{tab:recognition}): العتبة تعطي ”و“ على الأجزاء، والتثبيط المتبادل ”أو“ على المواضع، والقسمة الاستقلال عن السطوع، وقاعدة هيب مع أثر تعطي المعلّم الغائب. أخذ المهندسون من البصر تناوب الانتقائية واللاتغيّر وبنوا عليه الشبكات الالتفافية؛ أما الدماغ فلم يسلّمنا بعدُ أهم ما لديه: القدرة على التعلم من فيديو بلا أسماء، بلا طاقة تقريبًا.

\section{تمارين}

\begin{exercise}[\lvlA]
\label{ex:12:1}
\emph{التردد في مرحلة واحدة.} تدوم مرحلة خط الأنابيب $15\ms$، والعصبونات تعمل بتردد $20$ نبضة في الثانية، والضجيج مترابط بـ$c=0.1$~\eqref{eq:correlated}. ما الخطأ النسبي لتردد مجموعة من مئة عصبون، وما حده لمجموعة كبيرة بلا حد؟ كم مستوى من التردد يمكن تمييزه في المرحلة؟
\end{exercise}

\begin{exercise}[\lvlA]
\label{ex:12:2}
\emph{طاقة تعرّف واحد.} في مرور واحد خلال $150\ms$ تُصدر عشر مراحل من $10^7$ عصبون لكل منها نبضة واحدة على الأكثر، كلفتها $10^{-10}\,\text{J}$. (أ)~ما نسبة طاقة الدماغ كله خلال هذه الـ$150\ms$ التي يكلفها التعرّف؟ (ب)~كم مرة أكثر يستهلك مسرّع قدرته $300\,\text{W}$ يتعرف على الصورة في $10\ms$؟
\end{exercise}

\begin{exercise}[\lvlB]
\label{ex:12:3}
\emph{عتبات المرحلة $S$.} ”شبكية“ من $24$ نقطة، والشكلان $A=11011$ و$B=10101$، والنموذج~\eqref{eq:scstage} يساوي $+1$ على آحاد الشكل و$-1$ على أصفاره. (أ)~عند أي عتبات يتسامح عصبون $S$ مع نقطة مقلوبة واحدة، ويصمت على الشكل الآخر؟ (ب)~كم عصبون $S$ يلزم لعشر سمات $5\times5$ في صورة $100\times100$؟
\end{exercise}

\begin{exercise}[\lvlB]
\label{ex:12:4}
\emph{كم تدوم إحدى الصورتين.} في مخطط القفزات~\eqref{eq:halfcentre} عند $\tau\ll\tau_a$، ما دامت المجموعة~1 فائزة، يكون $r_2=0$ و$r_1=I-a_1$. (أ)~أوجد مدة الفوز عند $I=1$ و$\beta=2$ و$b=2$ و$\tau_a=0.4\,\text{s}$. (ب)~أي $\tau_a$ يعطي تبدّل الصورة كل ثلاث ثوانٍ؟
\end{exercise}

\begin{exercise}[\lvlB]
\label{ex:12:5}
\emph{نمذجة: ثمن اللاتغيّر.} بالدالة \foreignlanguage{english}{\texttt{two\_stage}} من \foreignlanguage{english}{\texttt{models/\allowbreak recognition\_models.py}} ($4000$ محاولة) أوجد عند أي احتمال لقلب النقطة $p$ يخطئ زوج المراحل $S$، $C$ عند $N=24$ في حالة من كل أربع. وكيف تتغير نسبة الأجوبة الصحيحة عند $p=0.1$ من $N=8$ إلى $N=192$؟
\end{exercise}

\begin{exercise}[\lvlB]
\label{ex:12:6}
\emph{نمذجة: نافذة الأثر.} عشرة تشغيلات لـ\foreignlanguage{english}{\texttt{trace\_learning}} من \foreignlanguage{english}{\texttt{models/\allowbreak recognition\_models.py}} بتوقف $0.3\,\text{s}$: أوجد عند أي قيم $\tau_{\text{tr}}$ بين $5\ms$ و$2\,\text{s}$ تتعلم كل التشغيلات اللاتغيّر بلا خطأ. ومن أين يأتي الحد الأعلى لهذه النافذة؟
\end{exercise}

\begin{exercise}[\lvlB]
\label{ex:12:7}
\emph{الحركة في أي مكان.} على نقاط متجاورة من صف بخطوة $0.1^\circ$ توجد كواشف اتجاه من الفصل~\ref{sec:gaba}: التثبيط من $A$ بتأخير $d$ يُطفئ الإثارة من $B$ إذا لم يتباعدا بأكثر من $5\ms$؛ وفوقها المرحلة $C$. ما التأخيرات اللازمة لكي تصمت $C$ عند الحركة المعاكسة بسرعة من $5$ إلى $20^\circ$ في الثانية؟
\end{exercise}

\begin{exercise}[\lvlC]
\label{ex:12:8}
\emph{تزييف البكسلات.} كم نقطة مقلوبة تلزم لتغيير جواب النموذج \foreignlanguage{english}{\texttt{two\_stage}} ($N=24$) على شكل نظيف؟ وكم تلزم للمصنِّف ”أكثر من $3.5$ آحاد في الدخل يعني $A$“، وهو أيضًا لا يتغير بالإزاحة؟ وما نسبة أجوبته الصحيحة عند $p=0.1$ مقابل $0.86$ لزوج المراحل؟
\end{exercise}
